预测MHC-I分子与肽段的结合，是肿瘤新抗原筛选、免疫治疗和免疫肽组学分析中的关键环节。随着人类蛋白质组中候选肽段规模持续膨胀，深度学习模型在实际部署时普遍面临高维肽段表征的缓存占用过大、存储成本高昂、I/O负担沉重以及在线查询效率低下等问题。为此，本研究以FennOmix.MHC为基础，聚焦大规模候选肽段场景，开展了模型轻量化与系统部署优化工作。

在复现模型并完成评估流程后，我们从排序能力、嵌入空间的结构保持度以及缓存开销等角度，比较了不同输出维度的方案。结果显示，64维嵌入能够在基本维持预测性能的同时，将肽段缓存从132 GB压缩至19 GB，在效果与存储成本之间取得了较好的平衡。接下来，我们对比了Linear和MLP两种降维结构，结合常规预测指标以及下游缩库辅助搜库的表现，选用了结构更简洁、部署代价更小的Linear作为默认降维方式。此外，通过对注意力头和前馈网络单元的重要性进行分析，确认模型存在冗余，适度剪枝后其预测表现依然保持稳定。

在系统优化方面，设计了一套基于INT8的缓存量化压缩策略，并配合自定义二进制缓存格式、哈希索引以及全量加载与延迟加载机制，构建了一条涵盖离线建库、索引定位、缓存读取、实时编码回退和在线匹配衔接的查询流水线。实验表明，该方案进一步缓解了缓存的存储和加载压力，提高了大规模候选肽段筛选的部署可行性。总体来看，本文从表征压缩、结构精简、缓存量化以及索引检索等多个维度对FennOmix.MHC进行了系统优化，为深度学习免疫肽结合预测模型的实际工程落地提供了一条可行的路径。